\section{The local rule}\label{sec:construction}

A \emph{rank} is the number of times a site has been encoded in a space--time block. The physical data track has rank zero. The rule acts at rank zero by the native recipe below. At higher ranks the recipe describes the computation performed on packets. It uses the same schedule and finite set of operations at every rank, with word width \(\ell_j=2^j\ell_0\). One tick and one site at rank \(j\) have physical sizes \(T^j\) and \(P^j\).

\subsection{States, bonds and structural updates}\label{sec:native}
The Purge, grades, provisional growth, Shrink and attack/withdrawal design adapts G\'acs \cite[Section~7(1--5), Section~9(M1--M12)]{gacs86}. The three evaluations and readout adapt his Section~5 and Section~9(M3--M5); the precise rule and guards used here are specified below.

Choose a sufficiently large integer \(w\), and set
\[
q=w^2,\quad A=w^3,\quad B=128q,\quad v=w^4,\quad
P_0=2^w,\quad P=5P_0,\quad d=4w+8 .
\]
The precise ordering of numerical choices is in Appendix~\ref{sec:constants}. The native state is either empty, \(\Dead\), or a live tuple with clock \(\tau\in[0,T)\), extended integer address \(\pi\in[-11P/10,21P/10)\), retained letter \(X\), bits \(\mathsf{Cn}\) and \(\mathsf{Temp}\), and the scratch registers below. A letter is either a distinct message-empty symbol or a finite tag with one word of width \(\ell_j\). Retained letters on live tuples are live. A live tuple can have arbitrary legal contents in the scratch registers.

The bond between an ordered left and right neighbor has grade zero if either is empty, if their clocks disagree, or if their residues modulo \(P\) do not increase by one modulo \(P\). Under the clock and residue conditions the bond has grade two when the integer address increases by one, grade one for a nonexact wrap from \(\pi>P/2\) to \(\pi<P/2\), and grade zero in other cases. The ordinary addresses belong to \([0,P)\). A true endpoint is address zero on the left or \(P-1\) on the right in the site's ordinary home. During the extended phase the true positions of a source are its own addresses zero and \(P-1\). The schedule depends on the old clock. Each tick uses the old row and updates simultaneously. Untargeted fields are retained. Surviving incumbents increment the clock modulo \(T\) and normalize their addresses modulo \(P\) in every ordinary update. Death clears the whole state.

Every recruitment is made into an old vacancy. A donor's clock must request a recruitment instruction, and its proposed address is its own plus the displacement to the vacancy. The proposed address must be legal. The neighbor opposite the donor must be empty or have the donor's clock. When both proposals qualify the left one is taken. A birth inherits the donor's next age and the proposed address, with ordinary normalization when required. Its retained letter is the corresponding live letter of the neutral packet encoding \(\Dead\). Its scratch is default, except that provisional births set \(\mathsf{Temp}=1\). Old incumbents do not accept proposals.

\begin{center}
\begin{tabular}{@{}p{.25\linewidth}p{.68\linewidth}@{}}
\toprule
Instruction & Structural effect \\
\midrule
Purge initialization & Set \(\mathsf{Cn}=1\). \\
Purge propagation & Set \(\mathsf{Cn}=0\) at a left bond below grade two or if the old left \(\mathsf{Cn}=0\). Retain it otherwise. \\
Purge kill & Kill when the old \(\mathsf{Cn}=0\) and the right bond has grade below two. \\
Provisional birth & Try both directions. During Grow use the extended legal range. In ordinary work demand an in-home address in \([0,P)\) on every proposal. \\
Provisional removal & Kill a marked incumbent with a grade-zero face unless that face is its own true outward face. \\
Grow attack & Try both directions, with proposed address outside the donor's ordinary home and in the legal extended range. \\
Grow withdrawal & Kill an exposed extended left tip (\(\pi<0\)) at a zero left bond, or an exposed extended right tip (\(\pi\ge P\)) at a zero right bond. \\
\(\operatorname{Shrink}(i)\) & Kill at a grade-zero face of depth greater than \(3iA\) from its corresponding true endpoint. \\
\bottomrule
\end{tabular}
\end{center}
At the end of Grow perform the final old-address withdrawal, then normalize every surviving address on entry to ordinary work. The next Purge initializes \(\mathsf{Cn}\) before a propagation or kill at the new address. A single exact source during Grow has a fixed value of physical child coordinate minus extended address, along with its clock. Normalization preserves residues and can merge labels that had distinct source identities.

\subsection{Clock}\label{sec:clock}

A handle has one Purge (initialization, \(q\) propagation ticks and \(q\) killing ticks) and one elementary instruction. Set \(H_0=2q+2\). Each Main cycle has \(2A+4\) handles. Its instructions, in order, are a reset of \(\mathsf{Temp}\) to zero, \(B\) simultaneous bilateral provisional attempts, \(B\) provisional removals, \(2A+2-2B\) idle instructions and the Main instruction. Its duration is \(d_0=(2A+4)H_0\). Write \(d_1=qd_0\).

Begin with \(\lfloor 7P/(5v)\rfloor\) Grow cycles, each of \(2v\) Main attacks followed by \(v\) Main withdrawals. Next perform \(\lfloor 3^w/q\rfloor\) calls of \(\operatorname{Shrink}(0)\), \(2P\) calls of \(\operatorname{Shrink}(1)\), and \(2P\) calls of \(\operatorname{Shrink}(2)\). Use one call in each stretch of \(q\) Main cycles, leaving the other \(q-1\) Main instructions idle. The end of Grow is \(s_0\), the end of the first Shrink stage is \(s_1\), and \(r_{\mathrm{last}}\) is the start of the final Grow withdrawal. The first Shrink stage accounts for more than \(0.85T\) of the full schedule at large \(w\).

Computation follows these three stages. Choose a service date \(s_c\) in early computation with more than \(w^6\) fine ticks before the first input transaction. Include three repetitions of the program below, separated by \(P\) idle Main cycles, and the verification and final ballot. The time for computation is a fixed constant times \(P\) plus a fixed polynomial in \(w\) in Main cycles. Within this allowance insert two no-birth retreat-through-padding intervals in the verifier phase separated by more than \(3w^6\) fine ticks, each containing more than \(100A+4(M+1)\) fine ticks of verification, with \(M=128d\). The second reaches the ballot, which is the last Main instruction of a normal cycle. The next period begins with a fresh Purge. The native handle cadence therefore continues across the vote. Verification runs on separate registers \emph{on fine ticks}, concurrently with Purge and the other scheduled operations. The no-birth portions of the normal cycles have more than enough fine ticks for the specified intervals. Use a fixed schedule satisfying these bounds, with the finite program below padded by idle instructions to the chosen time budget. This fixes \(T\) before any helper constants are selected. Since \(\operatorname{Shrink}(0)\) has duration of order \(3^w d_0\) and the other stages use \(O(Pq+\mathrm{poly}(w))\) Main cycles, the required fraction is achieved.

\subsection{Representation and evaluation}\label{sec:packets}

In addition to the structural fields, allocate three input and three output letters, two directional mailbox letters, a carrier \(Y\), two validation lanes each with a letter and a bit, a local invalid bit, a word-tape symbol, three external vector latches and a fixed tuple of parameter latch words. For the verifier allocate an active bit, two lane bits and two counters in \([0,M]\). Default scratch is blank, zero, headless, or message-empty in message fields. A diagnostic read of \(\Dead\) supplies canonical defaults and empty carriers. The word-tape symbol is the product of finitely many tape tracks, a finite control/head mark, and finitely many words. In particular the old head can carry two operands and one result. The head reads only its own old cell, writes the old cell's work tracks and requests a stationary move, a left move or a right move. For each cell resolve incoming requests in order stationary, from left, then from right. Impose work-interval barriers and an inert reject state on malformed instructions. These rules define an update even with multiple forged heads.

Serialize a rank-\((j+1)\) tuple in one rank-\(j\) packet of \(P_0\) letters. Use a live/dead discriminant and the canonical binary enumerations of all bounded fields in an \(O(w)\)-length area centered on the packet. Enumerate the registers in the declared order. Use separate scalar-bit tags with zero words and typed slots for the low and high word halves of each word field. Elsewhere require fixed live padding, including at both endpoints. Check tag and padding conventions at the packet level, reject out-of-range scalars and parse every legal tuple literally. In particular any legal mutable tag and payload, including headless registers with noncanonical content in the old state, is carried literally inside the typed word slots. A malformed packet parses as \(\Dead\). Encode \(\Dead\) using a playable all-live neutral packet with a dead discriminant and canonical body; the writer for an actual dead result instead emits all message-empty. For a live \(s\), its aligned encoding \(E_j(s)\) consists of five repeats of its serialized packet \(\delta_j(s)\): the site at block position \(a=cP_0+u\), with \(0\le c<5\) and \(0\le u<P_0\), has age zero, ordinary address \(a\), retained letter \(\delta_j(s)[u]\), \(\mathsf{Temp}=0\), \(\mathsf{Cn}=\mathbf 1_{\{a\ge q\}}\) and all scratch at default. These two structural bits are the values left at the ballot. The last complete Purge before it has a grade-one left wrap and grade-two interior bonds, so it leaves \(\mathsf{Cn}=1\) exactly at \(a\ge q\); the last Main cycle clears \(\mathsf{Temp}\) in the full home; and the ballot retains both bits. Define \(E_j(\Dead)=\Dead^P\), the all-empty block; it is distinct from the neutral packet that encodes \(\Dead\) above. Every field of each endpoint symbol of \(E_j(s)\), including the two structural bits, is independent of the particular live \(s\).

Number the fifths of a home \(0,1,2,3,4\). In each of three repetitions use the following ordered program. Reset the private output bank at its start.
\begin{enumerate}
\item For each supplier \(j'\in\{-1,0,1\}\) and fifth \(s\in\{1,2,3\}\), send its retained letter to bank \(s\) in fifth \(2+j'\) of the owner. For a displacement \(kP_0\), inject a fresh message on that lane at each designated source and message-empty everywhere else. Then shift \(|k|P_0\) times by unit nearest-neighbor shifts; each step copies the upstream mail only across a positive bond. Otherwise put message-empty. At receive, overwrite every addressed bank including message-empty. Zero displacement is local assignment, and births have empty mail.
\item Vote at each offset across the three homologous received letters. Any repeated letter wins. In a three-way tie choose the first live letter if at least two are live, and message-empty otherwise. In each supplied fifth independently check local tag, selector and padding constraints, and spread the OR of invalid bits for \(P_0\) steps within that fifth. Shift the two outer packets together with their invalid bits into the middle fifth. Overwrite the whole work tape, including heads, with the initialized interpreter and its three supplied input tracks.
\item At \(u=\lfloor P/2\rfloor,u+1,u+2\) load the actual read-only external words, and at \(u\) load the actual parameter. Gather these by two addressed nearest-neighbor transfers. The loads are unconditional on the old mutable tape and precede evaluation. The interpreter parses the supplied packets independently. It evaluates exactly one represented native update on the three parsed inputs using the actual parameter and vector words for any represented load operation. A halted true parameter returns \(\Dead\) before it evaluates any mutable instruction. Thus the guard applies equally to valid states with forged work registers.
\item Form the live result packet or put message-empty at every position for a dead result. Use a fine-site flag broadcast if needed before writing the full packet, and send the central fifth to each of the five output fifths with fresh injection, shifts and overwrite as in step 1. Maintain the completed bank while the other repetitions run.
\end{enumerate}

The interpreter may use copy, bitwise Boolean operations, modular addition, comparison, one-bit shifts, endpoint and zero tests, and population count on a word. Each operation on two half words uses a constant number of lower-width operations. For addition pass the low-half carry, for comparison compare the high halves first, and for shifts use the seam bit. The sum of the two half-population counts fits in one lower word for \(\ell_j\ge4\). Address and schedule arithmetic concerns \(O(w)\)-bit integers even when the simulated word is much larger. A fixed universal one-tape interpreter uses an initialization preamble and a quoted finite template \(p_{\mathrm{prog}}\) acting on the data \((p_{\mathrm{prog}},w)\). A represented initialization reconstructs the same preamble and supplied template from these data. When the represented state is a running head, evaluation performs its single elementary local instruction. This is a fixed finite tape algorithm. Compact-field scans and bounded arithmetic need a fixed polynomial budget \(w^{O(1)}<P_0\) at sufficiently large \(w\). Transfers, invalid broadcasts and fresh tape initialization use \(O(P_0)\) Main cycles. A budget of \(c_{\rm prog}P_0+w^{k_{\rm prog}}\) per repetition for fixed sufficiently large constants, plus the stated clock paddings, suffices. The table is specified by these finite algorithms and fixed timing constraints and can be enumerated by evaluating the resulting local map on every legal neighborhood.

\subsection{Verifier and ballot}\label{sec:verifier}

After readout reset the verification scratch and set the active bit to the majority of the three completed liveness bits at the same site. During each epoch of \(M+1\) fine ticks, inject each current active bit in its two directed lanes, clearing the counters, then shift the lane bits at unit speed for \(M\) ticks. Each shift uses a strict ordinary in-home bond with consecutive addresses and the same clock. A missing input is zero. Count positive arrivals in the two directions separately (saturating at \(M\)). A counter exceeding \(4d\) at the end of the epoch sets the local bit active. It remains active until death, reset or the ballot. Birth initializes all verifier fields to zero. At the final ballot an inactive site dies. An active site takes the payload vote by the same letter rule used at input, using a fixed live filler if that vote gives message-empty. The surviving site sets its age to zero, normalizes its address and clears every scratch register. The positive decision and the distributed payload are controlled separately in the induction.

\begin{lemma}[Clean full-row identity]\label{lem:clean}
On an aligned all-occupied row of encoded inputs with the stipulated read-only operands at the addressed loads throughout the three repetitions, one work period outputs exactly the aligned blocks \(E_j\) of the represented results, including \(E_j(\Dead)\). Every row at a date \(0\le u<T\) of the period is fully occupied, with the ordinary addresses of the inputs and the common clock \(u\); only the closing ballot at date \(T\) can make a result block empty.
\end{lemma}

\begin{proof}
The only weak bonds are aligned wraps a full home apart. A zero in the Purge propagation cannot reach the right end of a home in \(q\) ticks, so all Purges retain the complete row. Shrink, withdrawal and provisional removal find no zero face, and proposals cannot replace incumbents. Hence every row before the ballot is fully occupied. The input, work and readout transactions consequently compute and send the exact represented result in each repetition. In a home with a live result all the bits are initially active, and in a home with a dead result none can activate because the verifier's bonds are in-home. At the final vote the live result is encoded with its entire scratch reset, while the dead result becomes the all-empty block \(E_j(\Dead)=\Dead^P\), since every site of its home is inactive and dies at the ballot.
\end{proof}

\subsection{Read-only parameters}\label{sec:parameter}

For the fixed-rate family the physical literal is the exact initial configuration of \(U(n)\), hardwired into the native rule. The universal computation uses a one-ended tape with fixed bitplanes for symbols, state and head. At each promotion extend the word by a blank half and execute one true step. One step is implemented by a bounded list of masks, Boolean operations and one-bit shifts. The enlarged width accommodates the one newly visited site per step. If the literal for a transition is halted, its result is \(\Dead\) on every triple. Every substep reads only the old site and its two immediate neighbors: structural tests use those three old states, head requests and mail move one site, lane copies move one site, and the three addressed LOADs are gathered by two nearest-neighbor transfers. Local word operations read only local registers. The family rule is therefore radius one. The geometry and clock do not depend on the initial literal or word width. For the single rule, the fixed physical literal includes SEARCH, unary counter zero and the strict majority threshold for the number of external channels, with base width at least \(m_h+4\). SEARCH with three fresh words \(z_0,z_1,z_2\) compares
\[
\popcount\bigl((z_0\mathbin{\mathrm{xor}}z_1)\mathbin{\mathrm{and}}
                   (z_1\mathbin{\mathrm{xor}}z_2)\bigr)
\]
to that threshold. If a strict majority of channels make both flips, enter LOCKED with the true initial tape for the current counter index. Otherwise add one unary mark and use zero-extended \(z_0\) as the next external word. The external word after locking is zero and subsequent promotions advance the locked tape as in the family. The native step for any represented load takes these designated read-only operands. Mutated latch values remain ordinary mutable data.

\subsection{Physical embedding of the helper rule}\label{sec:embedding}

Use the auxiliary rule of Sections~\ref{sec:receiver}--\ref{sec:stripes}. Take a fixed stride \(D_{\rm lane}\) large enough to assign disjoint offsets in \([0,D_{\rm lane})\) to the data site and all private primitive roles, including every probe and individual replica. The data offset is zero. For every physical center \(x\), the declared register of replica offset \(a\) is held in the field for \(a\) in the physical symbol at \(x+a\). The symbol at each physical site has a field for each such offset, the data field, and a raw bit of deterministic proposal zero. Its deterministic local rule computes all its fields by translation of the respective logical proposals. Recursions at logical center \(x\) use logical neighbors \(x-D_{\rm lane},x,x+D_{\rm lane}\), with the appropriate finite offsets. The data update reads the exported vector at its own logical center. Auxiliary proposals never read data. Each physical update takes a single uniform whole-symbol replacement after computing the proposed tuple. In one residue lane the sites serving different primitive roles are distinct. Shifted lanes overlap in physical space, so no joint product across lanes is asserted. The physical alphabet is a finite product and the physical radius is finite and fixed. No geometric coordinate, noise rate, input index or externally supplied clock enters this translation-invariant update.
